\documentclass[12pt,letterpaper]{article}



%%
%% Required packages:
%%
\usepackage{etex}
\usepackage{amsmath}
\usepackage{amssymb}
\usepackage{amstext}
\usepackage{amsthm}
\usepackage{fancyhdr,fancybox}
\usepackage{mathrsfs} % need for \mathscr command
\usepackage{multicol}
\usepackage[normalem]{ulem}
%\usepackage{pictex,pictexplus}
% \usepackage{pictexwd}
\usepackage{rotating}
\usepackage{url}
\usepackage{xfrac}
\usepackage{booktabs}
\usepackage{color,soul}
\usepackage{comment}

\usepackage{hyperref}
\hypersetup{
    colorlinks=true,     % false: boxed links; true: colored links
    linkcolor=blue,      % color of internal links
    citecolor=blue,      % color of links to bibliography
    filecolor=blue,      % color of file links
    urlcolor=blue        % color of external links
}


\usepackage{tikz}
\usetikzlibrary{backgrounds,calc}

%%
%% Common formatting commands
%%
\setlength{\parindent}{0in}
\setlength{\textwidth}{7in}
\setlength{\evensidemargin}{-0.25in}
\setlength{\oddsidemargin}{-0.25in}
\setlength{\parskip}{.5\baselineskip}
\setlength{\topmargin}{-0.5in}
\setlength{\textheight}{9in}

%               Problem and Part
\newcounter{problemnumber}
\newcounter{partnumber}
\newcounter{subpartnumber}
\newcommand{\Problem}{%
  \refstepcounter{problemnumber}%
  \setcounter{partnumber}{0}%
  \setcounter{subpartnumber}{0}%
  \item[\makebox{\hfil\textbf{\theproblemnumber.}\hfil}]%
  }
\newcommand{\InProblem}[2][1.6]{%
  \refstepcounter{problemnumber}%
  \setcounter{partnumber}{0}%
  \setcounter{subpartnumber}{0}%
  \textbf{\theproblemnumber.}\ \ \parbox[t]{#1in}{#2}%
}
\newcommand{\InSmallProblem}[1]{\InProblem[1.05]{#1}}
\newcommand{\NoProblem}[1]{%
  \mbox{\hfil\textbf{#1}\hfil}%
  }
\newcommand{\UnProblem}{%
  \refstepcounter{problemnumber}%
  \setcounter{partnumber}{0}%
  \setcounter{subpartnumber}{0}%
  \mbox{\hfil\textbf{\theproblemnumber.}\hfil}%
  }
\newcommand{\InFlexProblem}[1]{%
  \refstepcounter{problemnumber}%
  \setcounter{partnumber}{0}%
  \setcounter{subpartnumber}{0}%
  \textbf{\theproblemnumber.}\ \ #1%
}
%%  Part:
\newcommand{\Part}{%
  \renewcommand{\thepartnumber}{(\alph{partnumber})}%
  \refstepcounter{partnumber}
  \setcounter{subpartnumber}{0}%
  \item[(\alph{partnumber})]%
}
\newcommand{\InPart}[2][1.75]{%
  \renewcommand{\thepartnumber}{(\alph{partnumber})}%
  \refstepcounter{partnumber}%
  \setcounter{subpartnumber}{0}%
  (\alph{partnumber})\ \ \parbox[t]{#1in}{#2}%
}
\newcommand{\UnPart}{%
  \refstepcounter{partnumber}%
  \renewcommand{\thepartnumber}{(\alph{partnumber})}%
  \setcounter{subpartnumber}{0}%
  (\alph{partnumber})%
}
\newcommand{\InLargePart}[1]{\InPart[3.05]{#1}}
\newcommand{\InFullPart}[1]{\InPart[6.2]{#1}}
\newcommand{\InSmallPart}[1]{\InPart[1.05]{#1}}
%% SubPart:
\newcommand{\SubPart}{%
  \refstepcounter{subpartnumber}%
  \item[(\roman{subpartnumber})]%
  }
\newcommand{\InSubPart}[2][2.25]{%
  \stepcounter{subpartnumber}%
  (\roman{subpartnumber})\ \ \parbox[t]{#1in}{#2}%
}
\newcommand{\InSmallSubPart}[1]{\InSubPart[1.5]{#1}}
\newcommand{\InTinySubPart}[1]{%
  \stepcounter{subpartnumber}%
  (\roman{subpartnumber})\ \ \makebox{#1}%
}
\newcommand{\InMySubPart}[2][2.25]{%
  \stepcounter{subpartnumber}%
  \parbox[t]{#1in}{\makebox[0.3in][l]{(\roman{subpartnumber})}\ \ {#2}}%
}
\newcommand{\TrueFalse}{\textbf{T}\ \ \textbf{F}\ \ }
\newcommand{\TFPart}[1]{% 
  \stepcounter{partnumber}%
  \setcounter{subpartnumber}{0}%
  \item[(\alph{partnumber})] \TrueFalse\parbox[t]{5.5in}{#1}}

\newcommand{\Points}[1]{(#1 points) \ }
\newcommand{\Point}[1]{(#1 point) \ }


%% For Answers:
\newcounter{answernumber}
\newcommand{\Answer}{%
  \refstepcounter{answernumber}%
  \setcounter{partnumber}{0}%
  \setcounter{subpartnumber}{0}%
  \item[\makebox{\hfil\textbf{\theanswernumber.}\hfil}]%
  }


% Setting up the headers for the first page
%\chead{} % will default to what is typed in the file.
\fancyhead[L]{\textsc{Math 22a}}
\fancyhead[R]{\textsc{Fall 2024}}
\fancyfoot[R]{
	\vspace*{\fill}\footnotesize\textcopyright{} \emph{Harvard Math 22a}	
}
\renewcommand{\headrulewidth}{0pt}

%\fancyhead[C]{}
%	\chead{}	
%	\lhead{}
%	\rhead{}
%	\cfoot{}


% Putting copyright on the footer of each page
%\fancypagestyle{secondpage}{
%	\fancyhead[C]{TESTing again} % change this to be blank
%		%\chead{}	
%	\fancyhead[L]{bears} % change this to be blank
%	\fancyhead[R]{dogs} % change this to be blank
%	\fancyfoot[R]{ %keeps the footer the same
%		\vspace*{\fill}\footnotesize\textcopyright{} \emph{Harvard Math 22a}	
%	}
%}
%\renewcommand{\headrulewidth}{0pt}
%\pagestyle{fancy}
\pagestyle{empty}


%\lhead{}
%\rhead{}
%\rfoot{\vspace*{\fill}\footnotesize\textcopyright{} \emph{Harvard Math 22a}}
%\renewcommand{\headrulewidth}{0pt}
%\pagestyle{fancy}




%%
%% Common shortcut commands:
%%
\newcommand{\ds}{\displaystyle}
\newcommand{\sgn}{\operatorname{sgn}}
\renewcommand{\Pr}{\operatorname{P}}
\newcommand{\CPr}[3][{}]{\Pr_{\text{#1}}\left(#2\,|\,#3\right)}

\newcommand{\C}{\mathbb{C}}
\newcommand{\R}{\mathbb{R}}
\newcommand{\Z}{\mathbb{Z}}
\newcommand{\N}{\mathbb{N}}
\newcommand{\Q}{\mathbb{Q}}
\newcommand{\D}{\mathbb{D}}

\newcommand{\rref}{\operatorname{rref}}
\newcommand{\im}{\operatorname{im}}
\newcommand{\rank}{\operatorname{rank}}
\newcommand{\diag}{\operatorname{diag}}
\newcommand{\Span}{\operatorname{span}}
%\renewcommand{\vec}[1]{\mathbf{#1}}
\newcommand{\charpoly}{\mathscr{P}}
\newcommand{\nicefrac}[2]{\sfrac{#1}{#2}} % using xfrac package
\newcommand{\topology}{\mathcal{T}}
\newcommand{\Inn}{\operatorname{Inn}}

\newcommand{\blankbox}[2]{\fbox{\rule{#1}{0in}\rule{0in}{#2}}}

\newenvironment{myindentpar}[1]%
 {\begin{list}{}%
         {\setlength{\leftmargin}{#1}
          \setlength{\rightmargin}{\leftmargin}}%
         \item[]%
 }
 {\end{list}}
\newtheorem{theorem}{Theorem}
\newtheorem*{NStheorem}{Theorem}
\newtheorem{cor}[theorem]{Corollary}
\newtheorem*{claim}{Claim}
\newtheorem{defn}{Definition}

\newenvironment{amatrix}[1]{%
  \left(\begin{array}{@{}*{#1}{c}|c@{}}
}{%
  \end{array}\right)
}

%--------grstep
% For denoting a Gauss' reduction step.
% Use as: \grstep{\rho_1+\rho_3} or \grstep[2\rho_5 \\ 3\rho_6]{\rho_1+\rho_3}
\newcommand{\grstep}[2][\relax]{%
   \ensuremath{\mathrel{
       {\mathop{\longrightarrow}\limits^{#2\mathstrut}_{
                                     \begin{subarray}{l} #1 \end{subarray}}}}}}
\newcommand{\swap}{\leftrightarrow}


\usetikzlibrary{arrows}
\usepackage{wrapfig}
\usepackage{amsfonts,amsmath}

\makeatletter
\renewcommand*\env@matrix[1][*\c@MaxMatrixCols c]{%
	\hskip -\arraycolsep
	\let\@ifnextchar\new@ifnextchar
	\array{#1}}
\makeatother

\def\env@matrix{\hskip -\arraycolsep
	\let\@ifnextchar\new@ifnextchar
	\array{*\c@MaxMatrixCols c}}

\newcommand{\norm}[1]{\left\lVert#1\right\rVert}

\chead{\Large\textbf{Dot Product and Geometry}}% 

\begin{document}
\thispagestyle{fancy}

Recall from last class:

\begin{itemize}
\item If $A$ is an $n \times n$ matrix, then $\lambda \in \mathbb{R}$ is an {\bf eigenvalue} of $A$ if there exists a non-zero vector $v$, called an {\bf eigenvector} corresponding to $\lambda$, such that $$Av=\lambda v.$$
\item If $\lambda$ is an eigenvalue of an $n \times n$ matrix $A$, then the set of all solutions $v$ to $Av = \lambda v$ is the {\bf eigenspace} of $A$ corresponding to $\lambda$.
\item {\bf Theorem 1:} The eigenvalues of a triangular matrix are the diagonal entries.
\item {\bf Theorem 2:} If $v_1, \dots, v_p$ are eigenvectors corresponding to distinct eigenvalues $\lambda_1, \dots, \lambda_p$ of an $n\times n$ matrix $A$, then $v_1, \dots, v_p$ are linearly independent.
\item If $A$ is an $n\times n$ matrix, then the characteristic polynomial of $A$, denoted by $p_A(\lambda)$, is given by $$p_A(\lambda)=\text{det}(A-\lambda I_n).$$
\item $\lambda$ is an eigenvalue of an $n\times n$ matrix $A$ if and only if $\lambda$ is a zero of the characteristic polynomial.
\item If $A$ and $B$ are similar matrices, then they have the same characteristic polynomials and hence the same eigenvalues (with the same algebraic multiplicities). 
\item An $n\times n$ matrix $A$ is {\bf diagonalizable} if $A$ is similar to a diagonal matrix.
\item {\bf Diagonalization Theorem:} An $n \times n$ matrix $A$ is diagonalizable if and only if $A$ has $n$ linearly independent eigenvectors. Moreover, $A=PDP^{-1}$ with $D$ diagonal if and only if the columns of $P$ are $n$ linearly independent eigenvectors of $A$. In this case the diagonal entires of $D$ are the corresponding eigenvalues of $A$.
\item {\bf Theorem 6:} An $n \times n$ matrix $A$ with $n$ distinct eigenvalues is diagonalizable.
\item 
{\bf Theorem 7:} Let $A$ be an $n\times n$ matrix with distinct eigenvalues $\lambda_1, \dots, \lambda_p$.
\begin{enumerate}
\item For $1\leq k \leq p$, the dimension of the eigenspace of $\lambda_k$ is at most the algebraic multiplicity of the eigenvalue $\lambda_k$.
\item The matrix $A$ is diagonalizable if and only if the sum of the dimensions of the eigenspaces is $n$ if and only if the characteristic polynomial factors into linear factors and the dimension for each eigenspace is exactly the algebraic multiplicity.
\item If $A$ is diagonalizable and $\mathcal{B}_k$ is a basis for the eigenspace corresponding to $\lambda_k$, then $\mathcal{B}_1, \dots, \mathcal{B}_p$ is an eigenbasis for $\mathbb{R}^n$.
\end{enumerate}
\end{itemize}

\newpage

\begin{defn} Let $u, v \in \mathbb{R}^n$, then the {\bf dot product} of $u$ and $v$ is given by $$u \cdot v= u^{T}v=u_1 v_1+\dots+ u_n v_n.$$
\end{defn}

{\bf Theorem 6.1:} Let $u, v, w \in \mathbb{R}^n$ and $c\in \mathbb{R}$.
\begin{enumerate}
\item $$ u \cdot v = v \cdot u.$$
\item $$(u+v)\cdot w= u \cdot w+ v \cdot w.$$
\item $$(cu)\cdot v = c(u\cdot v)= u \cdot (cv).$$
\item $u \cdot u \geq 0$ and $u \cdot u =0$ if and only if $u=0$.
\end{enumerate}

\begin{defn} The {\bf length} or {\bf norm} of the vector $v$ is a non-negative scalar, denoted by $\norm{v}$, and defined by $$\norm{v}= \sqrt{v\cdot v}.$$
\end{defn}

\begin{defn} We define the {\bf distance} between $u$ and $v$ by $$\text{dist}(u,v)=\norm{u-v}.$$
\end{defn}

\begin{defn} The vectors $u$ and $v$ are {\bf orthogonal} if $u\cdot v=0$.
\end{defn}

{\bf Theorem 6.2:} The vectors $u$ and $v$ are orthongal if and only if $$\norm{u}^2+\norm{v}^2=\norm{u+v}^2.$$

\begin{enumerate}

\Problem Prove the Theorem.

\vfill

\begin{defn} The set $$W^{\perp} =\{ z \in \mathbb{R}^n: z \cdot w=0 \text{ for all } w \in W\}$$ is called the {\bf orthogonal complement} of $W$.
\end{defn}

\newpage

{\bf Note:} $x \in W^{\perp}$ if and only if $x$ is orthogonal to each vector in a set that spans $W$. Also $W^{\perp}$ is a subspace of $\mathbb{R}^n$.

\vfill

{\bf Theorem 6.3} Let $A$ be an $m\times n$ matrix, then $$(\text{Row}(A))^{\perp} = \text{Null}(A) \text{ and } (\text{Col}(A))^{\perp}= \text{Null}(A^{T}).$$

\Problem Prove the Theorem.

\vfill

{\bf Cauchy-Schwarz Inequality:} For all ${\bf x}, {\bf y} \in {\mathbb R}^n$ we have
$$ |{\bf x} \cdot {\bf y}| \le \|{\bf x}\| \ \|{\bf y}\|. $$

\item Prove Cauchy-Schwarz.

\vfill

{\bf Note:} We can use the dot product to find the angles between vectors; namely $$u\cdot v = \norm{u} \norm{v} \cos \theta.$$

{\bf Triangle Inequality:} For all ${\bf x}, {\bf y} \in {\mathbb R}^n$  we have $\|{\bf x}+{\bf y}\| \le \|{\bf x}\| + \|{\bf y} \| $.

\item Prove the triangle inequality.

\vfill

\newpage

\begin{defn}
	A set of vectors $\{ u_1, \dots, u_p\}$ in $\mathbb{R}^n$ is {\bf orthogonal} if $u_k \cdot u_j=0$ for all $j\neq k$.
\end{defn}

{\bf Theorem 6.4:} If $S=\{ u_1, \dots, u_p\}$ is an orthogonal set of non-zero vectors in $\mathbb{R}^n$, then $S$ is linearly independent and hence a basis for $\text{Span} \; S$.

\Problem Prove the Theorem.

\vfill

\begin{defn} An {\bf orthogonal basis} for a subspace $W$ of $\mathbb{R}^n$ is a basis for $W$ that is also an orthogonal set.
\end{defn}


{\bf Theorem 6.5:} Let $\{u_1, \dots,u_p\}$ be an orthogonal basis for a subspace $W$ of $\mathbb{R}^n$. For each $y\in W$, then $$y=c_1 u_1+\dots+c_p u_p$$ where $$c_j=\frac{y\cdot u_j}{u_j \cdot u_j}.$$

\Problem Prove the Theorem.

\vfill

\Problem Let $u_1=\begin{bmatrix} 2\\ -3\\ \end{bmatrix}$ and $u_2= \begin{bmatrix} 6\\4\\ \end{bmatrix}$ and $x=\begin{bmatrix} 9\\ -7\\ \end{bmatrix}$. Find $x$ in the $\{u_1, u_2\}$-coordinates.



\vfill


\newpage




\end{enumerate}

\newpage
\chead{\Large\textbf{Dot Product -- Answers / Solutions}}%
\lhead{}
\rhead{}
\thispagestyle{fancy}

\begin{enumerate}
\Answer 

\begin{proof}
	Let $u, v \in \mathbb{R}^n$, then $$\norm{u+v}^2 =(u+v)\cdot (u+v)=u \cdot u +2 u \cdot v + v \cdot v =\norm{u}^2 +2 u\cdot v+ \norm{v}^2.$$ Thus $\norm{u+v}^2=\norm{u}^2+\norm{v}^2$ if and only if $u \cdot v =0$, which is equivalent to $u$ and $v$ are orthogonal.
\end{proof}

\Answer
\begin{proof}
	We first prove $(\text{Row}(A))^{\perp} = \text{Null}(A)$. Let $x \in (\text{Row}(A))^{\perp}$ and let $a_1, \dots, a_n$ be the rows of $A$. Since $x \in (\text{Row}(A))^{\perp}$, we know that $x \cdot a_k$ for all $1 \leq k \leq p$. Now we check if $x$ is in the null space of $A$: $$Ax = \begin{bmatrix} a_1 \cdot x\\ \vdots\\ a_n \cdot x\\ \end{bmatrix} =0.$$ Thus $x \in \text{Null}(A)$.
	
	Suppose $x \in \text{Null}(A)$, then $$0=Ax = \begin{bmatrix} a_1 \cdot x\\ \vdots\\ a_n \cdot x\\ \end{bmatrix}.$$ Thus $a_k \cdot x$ =0 for all $1\leq k \leq n$. Since $a_1, \dots, a_n$ is a spanning set for the row space of $A$, we know that $x$ is orhogonal to all of the row space. Thus $x \in \text{Row}(A)^\perp$.
	
	Now we prove $(\text{Col}(A))^{\perp}= \text{Null}(A^{T})$. Recall $\text{Col}(A)=\text{Row}(A^{T})$, and hence by the first part of the theorem, we get $$(\text{Col}(A))^{\perp}=(\text{Row}(A^{T}))^{\perp}= \text{Null}(A^{T}).$$
\end{proof}
\Answer For any scalar $t$ and any ${\bf x}, {\bf y}$ we have
$$ 0 \le \|{\bf x} + t {\bf y}\|^2 = \|{\bf x}\|^2 + 2 t {\bf x} \cdot {\bf y} + t^2 \|{\bf y}\|^2. \eqno (55) $$
The idea is to choose $t$ so as to minimize the expression on the right-hand side of (55).
Call the right-hand side $\phi(t)$. It is a quadratic polynomial in $t$, and it has its minimum when
$\phi'(t) = 0$. The corresponding value is when
$$ t = {- {\bf x} \cdot {\bf y} \over \|{\bf y}\|^2} .$$
Plugging in this value of $t$ and simplifying gives the Cauchy-Schwarz inequality. 
\Answer Since both sides of the inequality are non-negative, it holds if and only if it holds after squaring both sides. Doing so and cancelling terms
shows that it follows from Cauchy-Schwarz. 
\Answer
\begin{proof}
Suppose $c_1 u_1+\dots+	c_p u_p=0$. Taking the dot product with $u_k$, we get $$c_1 u_1 \cdot u_k + \dots +c_p u_p\cdot u_k=0,$$ and using orthogonality of $u_1, \dots, u_p$, we get $c_k u_k \cdot u_k=0$. Now since $u_k$ is non-zero, we get $c_k=0$. Thus $c_k=0$ for all $1\leq k \leq p$, and hence $u_1, \dots, u_p$ are linearly independent.
\end{proof}
\Answer
\begin{proof} Suppose $y= c_1 u_1+\dots c_pu_p$. Now take the dot product of both sides with $u_k$ to get $$y\cdot u_k = (c_1 u_1+\dots +c_p u_p)\cdot u_k.$$ Using the orthogonality of $u_1, \dots, u_p$, we can simplify the right hand side to get $$y \cdot u_k= c_k u_k \cdot u_k.$$ Therefore $$c_k=\frac{y \cdot u_k}{u_k\cdot u_k}.$$
\end{proof}
\Answer In order to find $x$ in the $u_1, u_2$ coordinates, we need to find scalars $c_1, c_2$ such that $$x=c_1 u_1+ c_2 u_2.$$ Notice that $u_1\cdot u_2=0$, and hence we can use the previous theorem to conclude $$c_1=\frac{x\cdot u_1}{u_1\cdot u_1}= \frac{39}{13}=3,$$ and $$c_2=\frac{x\cdot u_2}{u_2\cdot u_2}= \frac{26}{52}=\frac{1}{2}.$$ Thus $[x]_{\{u_1,u_1\}}=\begin{bmatrix} 3\\ \frac{1}{2}\\ \end{bmatrix}$.
\end{enumerate}  

\end{document}


%%% Local ables: 
%%% mode: latex
%%% TeX-master: t
%%% End: 
